\documentclass[12pt]{article}
\usepackage[T1]{fontenc}
\usepackage[utf8]{inputenc}
\usepackage{lmodern}
\usepackage{textcomp}
\usepackage{enumitem}
\usepackage{geometry}
\geometry{margin=1in}
\begin{document}
\begin{center}
Answer Key - History - K-12 Level Assessment 10 \
\small (Answers are ordered exactly as in the question paper.)
\end{center}

\section*{Multiple Choice}
\begin{enumerate}[label=\arabic*.]
\item \textbf{Answer:} D
\\ \textit{Options:} A) Mesopotamians.; B) Chavin.; C) Chinese.; D) Olmec.
\\ \textit{Note:} The Maya derived their writing and mathematics from the Olmec, as the Olmec civilization is recognized as one of the earliest Mesoamerican cultures to develop a system of hieroglyphic writing and a form of numerical notation that influenced subsequent civilizations, including the Maya.
\item \textbf{Answer:} C
\\ \textit{Options:} A) the height and size of its monumental architecture; B) the amount and frequency of violence; C) the degree of urban planning; D) the amount of surplus wheat that was produced for storage
\\ \textit{Note:} The Indus Valley civilization is considered unmatched in urban planning because it featured advanced city layouts with grid patterns, sophisticated drainage systems, and standardized brick sizes, which were highly developed compared to contemporaneous civilizations.
\item \textbf{Answer:} A
\\ \textit{Options:} A) dendrochronology.; B) coins recovered from ancient shipwrecks.; C) written records.; D) all of the above.
\\ \textit{Note:} Dendrochronology provides precise annual dating through tree ring analysis, allowing for more accurate calibration of radiocarbon dating results by establishing a chronological framework.
\item \textbf{Answer:} A
\\ \textit{Options:} A) political and military history.; B) philosophy and religion.; C) trade and economics.; D) astronomy and mathematics.
\\ \textit{Note:} The records left by the Maya primarily focus on political and military history because they document the achievements, conflicts, and lineage of rulers, contrasting with Mesopotamian writings that often emphasized trade and economic transactions.
\item \textbf{Answer:} A
\\ \textit{Options:} A) Europe; B) Africa; C) Asia; D) Australia
\\ \textit{Note:} Most of the Venus figurines, which are prehistoric statuettes associated with fertility and femininity, have been discovered in various locations across Europe, particularly in regions like the Upper Paleolithic sites of France and Austria.
\end{enumerate}

\section*{True / False}
\begin{enumerate}[label=\arabic*.]
\setcounter{enumi}{5}
\item \textbf{Answer:} False
\\ \textit{Options:} A) True; B) False
\\ \textit{Note:} The statement is false because, although Egyptian hieroglyphs have been largely deciphered, many aspects of the language and its usage remain unclear, making it complex and not fully understood.
\item \textbf{Answer:} False
\\ \textit{Options:} A) True; B) False
\\ \textit{Note:} The statement is false because the native peoples of the northwest coast primarily relied on fishing, hunting, and gathering rather than maize agriculture, which was more common in other regions of North America.
\item \textbf{Answer:} False
\\ \textit{Options:} A) True; B) False
\\ \textit{Note:} Mesopotamia is actually located between the Tigris and Euphrates rivers in the region known as the Fertile Crescent, not between the Amazon and Nile rivers.
\item \textbf{Answer:} True
\\ \textit{Options:} A) True; B) False
\\ \textit{Note:} Lewis Henry Morgan's stages of human culture, which include savagery, barbarism, and civilization, represent a unilineal evolutionary framework that posits all societies progress through these sequential stages in a linear fashion.
\item \textbf{Answer:} True
\\ \textit{Options:} A) True; B) False
\\ \textit{Note:} The florescence of Ancestral Puebloan culture, characterized by significant advancements in architecture, agriculture, and social organization, indeed occurred just after A.D. 1000 as they transitioned into a more complex society during the Pueblo II period.
\end{enumerate}

\section*{Long Form Response}
\begin{enumerate}[label=\arabic*.]
\setcounter{enumi}{10}
\item \textbf{Answer:} During the Upper Paleolithic period, items of personal adornment, such as jewelry made from shells, bones, and stones, reflect a growing awareness of individual identity, as discussed by Randall White. These adornments were not just practical but served as expressions of personal style and social status, indicating that individuals began to see themselves as unique beings within their communities. The presence of intricate designs and the use of specific materials suggest that people were beginning to value personal expression and identity more than before. This shift illustrates the importance of individuality in human societies, as these items allowed people to convey their personal stories and connections to others. Overall, personal adornments from this time mark a significant development in the understanding of self and identity in early human culture.
\\ \textit{Note:} correct because it highlights how the creation and use of personal adornments in the Upper Paleolithic period signify a shift towards individual identity, reflecting personal expression and social status, as emphasized by Randall White.
\item \textbf{Answer:} The intelligence of Homo erectus played a significant role in their ability to migrate more effectively than Homo habilis. Homo erectus had a larger brain and more advanced problem-solving skills, which allowed them to develop better tools and navigate diverse environments. This intelligence enabled them to plan their migrations, find food, and adapt to different climates more successfully than their predecessors. As a result, the successful migrations of Homo erectus contributed to the spread of human populations across continents, leading to greater genetic diversity and the eventual evolution of modern humans. This advancement in intelligence not only shaped the future of human migration but also laid the groundwork for more complex societies and cultures.
\\ \textit{Note:} correct because it identifies the larger brain size and advanced problem-solving skills of Homo erectus as key factors that enhanced their tool-making, adaptability, and planning abilities, facilitating more effective migrations compared to Homo habilis, which had significant implications for the evolution and spread of early humans.
\end{enumerate}

\vfill
\noindent\textit{End of Answer Key}
\end{document}